% ---- beta / M-c slope identifiability snippet (grounded in notebooks/_scratch_beta.py output) ----
% Intended for the Hierarchical SBI / discussion: explains why the 1D population marginals are
% recovered well while the M-c slope beta is not, across most experiments.

\paragraph{Why the slope $\beta$ is weakly identified.} A consistent feature of all three
hierarchical methods is that the one-dimensional population marginals---$\mu_M$, $\sigma_M$, $c_0$,
and the marginal concentration spread---are recovered accurately, while the mass--concentration
\emph{slope} $\beta$ is not: its posterior remains broad (typically $50$--$75\%$ of the prior width)
and its central value is biased toward steeper slopes. This is not a failure of any particular
method but a statement about what the data constrain. The slope $\beta = \mathrm{d}c/\mathrm{d}\log_{10}M$
can only be measured if the clusters span an appreciable \emph{range} in mass: with no lever arm in
$\log_{10}M$, the per-cluster $(M,c)$ pairs carry no information about how $c$ trends with $M$. The
controlling variable is therefore the population mass spread $\sigma_M$ itself. Across our
experiments the fractional width of the $\beta$ posterior is strongly anti-correlated with $\sigma_M$
(Pearson $r=-0.81$): in the high $\lambda$--$M$ scatter experiment, where $\sigma_M\approx0.54$
provides a wide mass baseline, $\beta$ is in fact recovered tightly and accurately
($-1.36$ vs.\ true $-1.33$, posterior width only $27\%$ of the prior), whereas in the experiments
with the fiducial narrow mass range ($\sigma_M\approx0.12$) the slope is essentially unconstrained.
Notably this is the same experiment whose 1D marginals look the most challenging---the slope is best
determined precisely where the mass distribution is broadest. We verified that this is a lever-arm
identifiability effect rather than the noise floor that limits $\sigma_c$: the $\beta$ posterior
width is uncorrelated with the per-bin measurement noise ($r\approx-0.01$), and the $\beta$ and
$\sigma_c$ posteriors are essentially uncorrelated ($r\approx-0.1$), so $\beta$ is not simply trading
off against the concentration scatter at fixed marginals (Figure~\ref{fig:beta_leverarm}). The
practical implication is that a precise $M$--$c$ \emph{slope} from stacked weak lensing requires a
wide intrinsic mass baseline, which a single richness bin does not provide.

We confirm this directly with a joint fit across all seven richness bins. Pooling the bins spans
$\sim$2.1\,dex in $\log_{10}M$ (versus $\sim$0.12 within a single bin), an order-of-magnitude wider
lever arm. We fit one shared $M$--$c$ relation across the bins while giving each bin its own
mass-function$\times$selection population (the same machinery as the ablation above,
Figure~\ref{fig:massfn}), so that the per-bin selection is handled self-consistently and $\beta$ is constrained by how the concentration
tracks mass across the full baseline. The joint fit recovers $\beta=-0.86\pm0.04$ against a true
cross-bin slope of $-0.97$---a posterior $7\times$ tighter than the single-bin value
($\beta=-1.28\pm0.31$). The slope, unconstrained within any one bin, is thus well determined once the
mass baseline is wide, exactly as the lever-arm picture predicts (Figure~\ref{fig:multibin_beta}); we
recommend the multi-bin joint fit whenever the $M$--$c$ slope itself is of interest.

\begin{figure*}[t]
\centering
\includegraphics[width=0.92\textwidth]{figures/multibin_beta.png}
\caption{Recovering the $M$--$c$ slope $\beta$ with a multi-bin joint fit. \emph{Left:} the pooled
$(M,c)$ population of all seven richness bins (colored by bin) spans $\sim$2.1\,dex in $\log_{10}M$,
an order-of-magnitude wider mass baseline than any single bin; the joint-fit relation (red, band =
$1\sigma$) tracks the true cross-bin slope (black dashed). \emph{Right:} the $\beta$ posterior is
$7\times$ tighter for the multi-bin fit (red, $-0.86\pm0.04$) than for a single bin (purple,
$-1.28\pm0.31$), and brackets the truth (dashed). The slope that is unconstrained within one richness
bin becomes well determined once the bins are fit jointly.\label{fig:multibin_beta}}
\end{figure*}

\begin{figure}[t]
\centering
\includegraphics[width=\columnwidth]{figures/beta_leverarm.png}
\caption{The $M$--$c$ slope $\beta$ is identifiable only with a mass lever arm. \emph{Left:} the
fractional width of the $\beta$ posterior (relative to its prior) versus the true population mass
spread $\sigma_M$, for the Hierarchical SBI and Hybrid SBI--HMC; the two are strongly anti-correlated
($r=-0.81$), so $\beta$ is constrained only where the clusters span an appreciable mass range.
\emph{Right:} recovered versus true $\beta$ (marker size $\propto\sigma_M$); the wide-mass-range
experiment lands on the one-to-one line while the narrow-mass experiments scatter and rail toward the
prior.\label{fig:beta_leverarm}}
\end{figure}

\paragraph{A realistic-survey demonstration.} To test this on a survey-realistic dataset rather than
the controlled equal-count bins, we assemble the full \citet{mcclintock18} sample geometry: the four
richness bins above their calibration floor ($\lambda\in[20,30),[30,45),[45,60),[60,\infty)$) crossed
with three redshift bins ($z\in[0.2,0.35),[0.35,0.5),[0.5,0.65)$), using the actual per-cell cluster
counts from that analysis (6504 clusters in total, ranging from 1612 in the most populated cell to 91
in the sparsest). We fit one shared $M$--$c$ relation across all twelve cells, now including a
redshift-evolution term, $c = c_0 + \beta(\log_{10}M - 14.3) + \gamma(z-0.35) + \mathcal{N}(0,\sigma_c)$,
with each cell carrying its own mass-function$\times$selection mass population at its own redshift
(and a redshift-aware forward model). We run this with both hierarchical HMC and the Hybrid SBI--HMC.
Both recover the slope that is inaccessible in any single bin (Figure~\ref{fig:fulldataset}): the
Hybrid gives $\beta=-0.70\pm0.07$ and HMC $\beta=-0.56\pm0.06$ against a true $-0.68$, posteriors
$\sim5\times$ tighter than the single-bin value and no longer prior-railed. The redshift evolution,
which only this multi-redshift dataset can probe, is also detected ($\gamma=-2.09\pm0.10$ from HMC
versus a true $-2.07$). The two methods differ mildly on the harder second-order parameters---HMC
slightly underestimates $\beta$ and $\sigma_c$ (a residual interaction with the per-cell
mass-function prior), while the Hybrid's $\gamma$ is pulled somewhat steep (a calibration drift of the
redshift-conditioned per-cluster network at the redshift extremes)---but both deliver the central
result: with a realistic survey's mass and redshift baseline, the $M$--$c$ slope and its redshift
evolution are well constrained, where a single richness bin determines neither.

\begin{figure*}[t]
\centering
\includegraphics[width=0.95\textwidth]{figures/fulldataset_recovery.png}
\caption{Full-dataset recovery on the realistic \citet{mcclintock18} $4\times3$ (richness $\times$
redshift) grid, 6504 clusters. \emph{Left:} the pooled $(M,c)$ population (points colored by redshift)
with the recovered $M$--$c$ relation at each redshift; the relation shifts downward with $z$ (the
recovered evolution $\gamma$) and tilts with mass (the slope $\beta$). \emph{Right:} recovered $\beta$
and $\gamma$ for hierarchical HMC and Hybrid SBI--HMC against truth; both recover the slope that is
unconstrained in any single richness bin ($\beta=-1.28\pm0.31$, biased) and detect the
concentration--redshift evolution.\label{fig:fulldataset}}
\end{figure*}
